\documentclass[12pt]{article}
\usepackage[T2A]{fontenc} 
\usepackage[utf8]{inputenc} 
\usepackage[russian]{babel} 
\usepackage{amsmath, amssymb, amsthm}
\usepackage{geometry}
\geometry{a4paper, margin=2cm}
\usepackage{enumitem}

\begin{document}
	
	\title{}
	\author{}
	\date{}
	\maketitle
	
	\textbf{Основная часть.}
	
	\bigskip
	\noindent\textbf{Задача 1.}
	
	\medskip
	\textbf{1)} \quad
	\[
	X = \{ x=(x_1,x_2)\in\mathbb{R}^2 : x_1^2+(x_2-1)^2\le 1 \}.
	\]
	Заметим, что $X$ --- замкнутый круг единичного радиуса с центром в точке $(0,1)$. Его опорная функция имеет вид
	\[
	h_X(y)=\max_{x\in X}\langle y,x\rangle=\langle y,(0,1)\rangle + \|y\|, 
	\]
	так как для шара $B(a,1)$ верно
	\[
	\max_{x\in B(a,1)}\langle y,x\rangle=\langle y,a\rangle+\|y\|.
	\]
	Следовательно, по определению полярного множества,
	\[
	X^*=\{y\in\mathbb{R}^2: y_2+\|y\|\le 1\}.
	\]
	
	\medskip
	\textbf{2)} \quad
	\[
	X=\{ x=(x_1,x_2)\in\mathbb{R}^2: x_2\ge \tfrac{x_1^2}{2} \}.
	\]
	Заметим, что множество $X$ неограничено. Рассмотрим опорную функцию:
	\[
	h_X(y)=\sup\{ y_1x_1+y_2x_2 : x_2\ge \tfrac{x_1^2}{2}\}.
	\]
	Если $y_2>0$, то, выбирая $x_1$ сколь угодно большим (и подставляя $x_2=\tfrac{x_1^2}{2}$), можно сделать значение
	\[
	y_1x_1+y_2\frac{x_1^2}{2}=\frac{y_2}{2}x_1^2+y_1x_1
	\]
	сколь угодно большим, то есть $h_X(y)=+\infty$. Для конечности опорной функции необходимо, чтобы $y_2\le 0$. При $y_2<0$ функция
	\[
	\varphi(x_1)=\frac{y_2}{2}x_1^2+y_1x_1
	\]
	является квадратичной и достигает максимума при $x_1=-\frac{y_1}{y_2}$. Подставляя, получаем:
	\[
	\varphi\Bigl(-\frac{y_1}{y_2}\Bigr)=\frac{y_2}{2}\left(\frac{y_1^2}{y_2^2}\right)-\frac{y_1^2}{y_2}=-\frac{y_1^2}{2y_2}.
	\]
	Требование $h_X(y)=\sup_{x_1\in\mathbb{R}}\varphi(x_1)\le 1$ эквивалентно неравенству
	\[
	-\frac{y_1^2}{2y_2}\le 1.
	\]
	Учитывая, что $y_2\le 0$, это преобразуется к условию
	\[
	y_1^2\le -2y_2.
	\]
	При $y_2=0$ для конечности супремума необходимо, чтобы $y_1=0$. Таким образом,
	\[
	X^*=\{(y_1,y_2)\in\mathbb{R}^2: y_2\le 0 \text{ и } y_1^2\le -2y_2\}.
	\]
	
	\medskip
	\textbf{3)} \quad
	\[
	X=\{ x\in\mathbb{R}^d: \|Ax\|\le 1 \}, \quad A\in S_d^{++},
	\]
	где $S_d^{++}$ --- множество симметричных положительно определённых матриц.
	\\
	Для данного $A$ оператор $x\mapsto Ax$ инвертируем, и опорная функция множества $X$ равна
	\[
	h_X(y)=\sup\{ \langle y,x \rangle: \|Ax\|\le 1\} = \|A^{-1}y\|,
	\]
	так как норма двойственная к $\|\cdot\|$ совпадает с ней. Следовательно,
	\[
	X^*=\{y\in\mathbb{R}^d: \|A^{-1}y\|\le 1\}.
	\]
	
	\bigskip

	Заметим, что множество $X$ является аффинным гиперплоскостью, не содержащей нуля. Тем не менее, по определению
	\[
	X^*=\{ y\in\mathbb{R}^d: \sup_{x\in X}\langle y,x\rangle\le 1\}.
	\]
	Заметим, что для любого $x\in X$ можно представить его в виде
	\[
	x=x_0+z, \quad \text{где } x_0 \text{ --- фиксированное решение } \sum_{i=1}^{d} (x_0)_i=1, \quad z\in \ker\Bigl(\sum_{i=1}^{d} x_i\Bigr).
	\]
	Если найдётся такой $y$, для которого существует $z$ с $\langle y,z\rangle\neq 0$, то за счёт неограниченности подпространства $\ker$ супремум будет равен $+\infty$. Следовательно, чтобы супремум был конечным, необходимо, чтобы
	\[
	\langle y,z\rangle=0 \quad \forall z\in \ker\Bigl(\sum_{i=1}^{d} x_i\Bigr),
	\]
	то есть $y$ должно быть коллинеарно вектору $(1,1,\dots,1)$. Пусть $y=\lambda (1,1,\dots,1)$. Тогда для любого $x\in X$ имеем
	\[
	\langle y,x\rangle=\lambda\sum_{i=1}^d x_i=\lambda.
	\]
	Требование $\sup_{x\in X}\langle y,x\rangle=\lambda\le 1$ означает, что
	\[
	\lambda\le 1.
	\]
	Таким образом,
	\[
	X^*=\{ \lambda (1,1,\dots,1) : \lambda\in\mathbb{R},\ \lambda\le 1 \}.
	\]
	
	Теперь найдём $X^{**}$ по определению:
	\[
	X^{**}=\{ x\in\mathbb{R}^d: \sup_{y\in X^*} \langle y,x\rangle\le 1 \}.
	\]
	Для произвольного $x\in\mathbb{R}^d$ и любого $y=\lambda(1,\dots,1)$ из $X^*$ получаем
	\[
	\langle y,x\rangle=\lambda\sum_{i=1}^d x_i.
	\]
	Заметим, что если $S=\sum_{i=1}^d x_i>0$, то супремум по $\lambda\le 1$ равен $S$. Если же $S<0$, то, поскольку $\lambda$ может принимать сколь угодно малые значения (отрицательные без нижней границы), супремум равен $+\infty$. Поэтому условие $\sup_{y\in X^*} \langle y,x\rangle\le 1$ возможно только если
	\[
	S\ge 0 \quad \text{и} \quad S\le 1.
	\]
	То есть
	\[
	X^{**}=\{ x\in\mathbb{R}^d: 0\le \sum_{i=1}^d x_i\le 1 \}.
	\]
	
	\vspace{1em}
	\textbf{Дополнительная часть.}
	
	\bigskip

	\textbf{1)Доказательство:}  
	Определим сопряжённое множество (в смысле полярного множества) для множества матриц по стандартной формуле:
	\[
	(A_d)^*=\{ Y\in\mathbb{R}^{d\times d} : \sup_{X\in A_d} \operatorname{tr}(YX) \le 1 \}.
	\]
	Пусть $Y\in S_d$, то есть $Y^T=Y$. Тогда для любой антисимметричной матрицы $X$ выполнено
	\[
	\operatorname{tr}(YX)=\operatorname{tr}(YX)^T=\operatorname{tr}(X^TY^T)=\operatorname{tr}((-X)Y)=-\operatorname{tr}(YX).
	\]
	Откуда следует, что $\operatorname{tr}(YX)=0$, и, следовательно, условие $\operatorname{tr}(YX)\le 1$ выполняется. Таким образом, $S_d\subset (A_d)^*$.
	
	Обратно, если $Y\not\in S_d$, то его симметричная и антисимметричная части
	\[
	Y=\frac{Y+Y^T}{2}+\frac{Y-Y^T}{2}=: Y_s + Y_a
	\]
	при этом $Y_a\neq 0$. Но поскольку для любой антисимметричной матрицы $X$ можно добиться произвольного значения $\operatorname{tr}(Y_aX)$ (используя масштабирование), супремум $\sup_{X\in A_d}\operatorname{tr}(YX)$ оказывается неограниченным. Следовательно, $Y\not\in (A_d)^*$. Получаем, что $(A_d)^*=S_d$.
	

	\medskip
	\textbf{2)Доказательство (эскиз):}  
	Пусть $(x,\mu)\in K_p$, то есть $\|x\|_p\le \mu$ с $\mu\ge 0$, и пусть $(y,\nu)\in\mathbb{R}^{d+1}$. Тогда по неравенству Гёльдера имеем
	\[
	\langle x,y\rangle\le \|x\|_p \|y\|_{p^*}\le \mu \|y\|_{p^*}.
	\]
	Отсюда для скалярного произведения получаем
	\[
	\langle (x,\mu), (y,\nu) \rangle = \langle x,y\rangle+\mu\nu \le \mu(\|y\|_{p^*}+\nu).
	\]
	Поскольку $K_p$ --- конус (то есть если $(x,\mu)\in K_p$, то и $(tx,t\mu)\in K_p$ для любого $t>0$), условие
	\[
	\sup_{(x,\mu)\in K_p} \langle (x,\mu), (y,\nu) \rangle \le 1
	\]
	может быть удовлетворено лишь при выполнении соотношения
	\[
	\|y\|_{p^*}+\nu\le 0.
	\]
	При соответствующем выборе нормировки (обычно поляром конуса называют множество всех $(y,\nu)$, для которых $\langle (x,\mu), (y,\nu) \rangle\le 0$ для всех $(x,\mu)\in K_p$), с учётом однородности, можно показать, что
	\[
	(K_p)^* = \{ (y,\nu)\in\mathbb{R}^{d+1} : \|y\|_{p^*}\le \nu \} =: K_{p^*}.
	\]
	Аналогичным образом доказывается, что $(K_{p^*})^* = K_p$.
	
	В доказательстве ключевую роль играет неравенство Гёльдера, которое обеспечивает точную взаимосвязь между нормами $\|\cdot\|_p$ и $\|\cdot\|_{p^*}$, а также однородность рассматриваемых конусов. В случае $p=1$ и $p^*=\infty$ рассуждения проводятся аналогично с использованием соответствующих норм.
	
	
\end{document}
